\documentclass[10pt,border=1pt]{standalone}
\usepackage{amsmath,amssymb,booktabs}
\usepackage{atlas-figs}
\begin{document}%
\begin{tabular}{@{}c@{}}\begin{tikzpicture}[bf,x=1cm,y=1cm]
\draw[black!25,line width=.3pt] (0.000,.4) -- (0.000,-8.83);\node[font=\footnotesize,anchor=south] at (0.000,.4) {$8$};
\draw[black!25,line width=.3pt] (1.190,.4) -- (1.190,-8.83);\node[font=\footnotesize,anchor=south] at (1.190,.4) {$12$};
\draw[black!25,line width=.3pt] (2.212,.4) -- (2.212,-8.83);\node[font=\footnotesize,anchor=south] at (2.212,.4) {$17$};
\draw[black!25,line width=.3pt] (3.223,.4) -- (3.223,-8.83);\node[font=\footnotesize,anchor=south] at (3.223,.4) {$24$};
\draw[black!25,line width=.3pt] (4.158,.4) -- (4.158,-8.83);\node[font=\footnotesize,anchor=south] at (4.158,.4) {$33$};
\draw[black!25,line width=.3pt] (5.377,.4) -- (5.377,-8.83);\node[font=\footnotesize,anchor=south] at (5.377,.4) {$50$};
\draw[black!25,line width=.3pt] (6.566,.4) -- (6.566,-8.83);\node[font=\footnotesize,anchor=south] at (6.566,.4) {$75$};
\draw[black!25,line width=.3pt] (7.410,.4) -- (7.410,-8.83);\node[font=\footnotesize,anchor=south] at (7.410,.4) {$100$};
\draw[black!25,line width=.3pt] (8.600,.4) -- (8.600,-8.83);\node[font=\footnotesize,anchor=south] at (8.600,.4) {$150$};
\node[font=\footnotesize,anchor=south west,inner sep=1pt] at (0,.95) {number of vertices of $T$};
\node[font=\footnotesize\bfseries,anchor=east,inner sep=1pt] at (-.2,.2) {open, and tested up to};
\draw[line width=1.5pt,dash pattern=on 3pt off 2pt] (2.968,-0.68) -- (7.410,-0.68);\node[font=\footnotesize,fill=white,inner sep=1pt,anchor=west] at (7.460,-0.68) {$100$};
\draw[line width=.7pt,fill=white] (0,-0.79) rectangle (2.968,-0.5700000000000001);
\fill[black!80] (0,-0.79) rectangle (2.968,-0.5700000000000001);
\node[font=\footnotesize,fill=white,inner sep=1pt,anchor=south west] at (3.008,-0.6000000000000001) {$22$};
\node[font=\footnotesize,anchor=east,text width=6.1cm,align=right,inner sep=1pt] at (-.2,-0.68) {at least $2^m$ good states, $m$ the fewest cycles of a cover: every root};
\draw[line width=.7pt,fill=white] (0,-1.4700000000000002) rectangle (3.223,-1.25);
\fill[black!80] (0,-1.4700000000000002) rectangle (2.034,-1.25);
\node[font=\footnotesize,fill=white,inner sep=1pt,anchor=south west] at (3.263,-1.28) {$16, 24$};
\node[font=\footnotesize,anchor=east,text width=6.1cm,align=right,inner sep=1pt] at (-.2,-1.36) {the same bound, on triangulations without separating triangles};
\draw[line width=1.5pt,dash pattern=on 3pt off 2pt] (2.034,-2.04) -- (8.600,-2.04);\node[font=\footnotesize,fill=white,inner sep=1pt,anchor=west] at (8.650,-2.04) {$150$};
\draw[line width=.7pt,fill=white] (0,-2.15) rectangle (2.034,-1.93);
\node[font=\footnotesize,fill=white,inner sep=1pt,anchor=south west] at (2.074,-1.96) {$16$};
\node[font=\footnotesize,anchor=east,text width=6.1cm,align=right,inner sep=1pt] at (-.2,-2.04) {greedy descent in the lattice reaches a good state};
\draw[line width=1.5pt,dash pattern=on 3pt off 2pt] (2.212,-2.72) -- (7.945,-2.72);\node[font=\footnotesize,fill=white,inner sep=1pt,anchor=west] at (7.995,-2.72) {$120$};
\draw[line width=.7pt,fill=white] (0,-2.83) rectangle (2.212,-2.6100000000000003);
\node[font=\footnotesize,fill=white,inner sep=1pt,anchor=south west] at (2.252,-2.64) {$17$};
\node[font=\footnotesize,anchor=east,text width=6.1cm,align=right,inner sep=1pt] at (-.2,-2.72) {exchange descent on rainbow trees reaches a good state};
\draw[line width=.7pt,fill=white] (0,-3.5100000000000002) rectangle (2.538,-3.2900000000000005);
\fill[black!80] (0,-3.5100000000000002) rectangle (2.212,-3.2900000000000005);
\node[font=\footnotesize,fill=white,inner sep=1pt,anchor=south west] at (2.578,-3.3200000000000003) {$17, 19$};
\node[font=\footnotesize,anchor=east,text width=6.1cm,align=right,inner sep=1pt] at (-.2,-3.4000000000000004) {some contraction keeps a pair};
\draw[line width=.7pt,fill=white] (0,-4.19) rectangle (2.379,-3.97);
\fill[black!80] (0,-4.19) rectangle (2.212,-3.97);
\node[font=\footnotesize,fill=white,inner sep=1pt,anchor=south west] at (2.419,-4.0) {$17, 18$};
\node[font=\footnotesize,anchor=east,text width=6.1cm,align=right,inner sep=1pt] at (-.2,-4.08) {a contraction never has more good states than the original};
\node[font=\footnotesize\bfseries,anchor=east,inner sep=1pt] at (-.2,-4.8) {refuted, counterexample known at};
\draw[line width=.4pt,black!55] (0,-5.3) -- (1.190,-5.3);\draw[line width=1.6pt] (1.090,-5.3999999999999995) -- (1.290,-5.2) (1.090,-5.2) -- (1.290,-5.3999999999999995);
\node[font=\footnotesize,inner sep=1pt,anchor=west] at (1.440,-5.3) {$12$};
\node[font=\footnotesize,anchor=east,text width=6.1cm,align=right,inner sep=1pt] at (-.2,-5.3) {the highest state with a partner is good};
\draw[line width=.4pt,black!55] (0,-5.8) -- (2.379,-5.8);\draw[line width=1.6pt] (2.279,-5.8999999999999995) -- (2.479,-5.7) (2.279,-5.7) -- (2.479,-5.8999999999999995);
\draw[line width=.8pt] (2.379,-5.8) circle[radius=.19];
\node[font=\footnotesize,inner sep=1pt,anchor=west] at (2.629,-5.8) {$18$};
\node[font=\footnotesize,anchor=east,text width=6.1cm,align=right,inner sep=1pt] at (-.2,-5.8) {an extreme cover has a good orientation};
\draw[line width=.4pt,black!55] (0,-6.3) -- (2.379,-6.3);\draw[line width=1.6pt] (2.279,-6.3999999999999995) -- (2.479,-6.2) (2.279,-6.2) -- (2.479,-6.3999999999999995);
\draw[line width=.8pt] (2.379,-6.3) circle[radius=.19];
\node[font=\footnotesize,inner sep=1pt,anchor=west] at (2.629,-6.3) {$18$};
\node[font=\footnotesize,anchor=east,text width=6.1cm,align=right,inner sep=1pt] at (-.2,-6.3) {some good state of a contraction lifts directly};
\draw[line width=.4pt,black!55] (0,-6.8) -- (3.223,-6.8);\draw[line width=1.6pt] (3.123,-6.8999999999999995) -- (3.323,-6.7) (3.123,-6.7) -- (3.323,-6.8999999999999995);
\node[font=\footnotesize,inner sep=1pt,anchor=west] at (3.473,-6.8) {$24$};
\node[font=\footnotesize,anchor=east,text width=6.1cm,align=right,inner sep=1pt] at (-.2,-6.8) {a cover one step from an extreme has one};
\draw[line width=.4pt,black!55] (0,-7.3) -- (4.158,-7.3);\draw[line width=1.6pt] (4.058,-7.3999999999999995) -- (4.258,-7.2) (4.058,-7.2) -- (4.258,-7.3999999999999995);
\node[font=\footnotesize,inner sep=1pt,anchor=west] at (4.408,-7.3) {$33$};
\node[font=\footnotesize,anchor=east,text width=6.1cm,align=right,inner sep=1pt] at (-.2,-7.3) {an extreme or irreducible cover has one};
\draw[line width=.4pt,black!55] (0,-7.8) -- (4.158,-7.8);\draw[line width=1.6pt] (4.058,-7.8999999999999995) -- (4.258,-7.7) (4.058,-7.7) -- (4.258,-7.8999999999999995);
\node[font=\footnotesize,inner sep=1pt,anchor=west] at (4.408,-7.8) {$33$};
\node[font=\footnotesize,anchor=east,text width=6.1cm,align=right,inner sep=1pt] at (-.2,-7.8) {every maximal chain of covers has one};
\draw[line width=.4pt,black!55] (0,-8.3) -- (4.158,-8.3);\draw[line width=1.6pt] (4.058,-8.4) -- (4.258,-8.200000000000001) (4.058,-8.200000000000001) -- (4.258,-8.4);
\node[font=\footnotesize,inner sep=1pt,anchor=west] at (4.408,-8.3) {$33$};
\node[font=\footnotesize,anchor=east,text width=6.1cm,align=right,inner sep=1pt] at (-.2,-8.3) {at least half of all covers have one};
\draw[line width=.4pt,black!55] (0,-8.8) -- (4.572,-8.8);\draw[line width=1.6pt] (4.472,-8.9) -- (4.672,-8.700000000000001) (4.472,-8.700000000000001) -- (4.672,-8.9);
\node[font=\footnotesize,inner sep=1pt,anchor=west] at (4.822,-8.8) {$38$};
\node[font=\footnotesize,anchor=east,text width=6.1cm,align=right,inner sep=1pt] at (-.2,-8.8) {the evict walk never repeats};
\fill[black!80] (-6.2,-9.66) rectangle (-5.4,-9.440000000000001);\node[right,font=\footnotesize] at (-5.4,-9.55) {every case, two programs};
\draw[line width=.7pt,fill=white] (-1.7000000000000002,-9.66) rectangle (-0.9000000000000004,-9.440000000000001);\node[right,font=\footnotesize] at (-0.9000000000000004,-9.55) {every case, one program};
\draw[line width=1.5pt,dash pattern=on 3pt off 2pt] (2.8,-9.55) -- (3.6000000000000005,-9.55) node[right,font=\footnotesize] {random samples};
\draw[line width=1.6pt] (6.2,-9.65) -- (6.3999999999999995,-9.450000000000001) (6.2,-9.450000000000001) -- (6.3999999999999995,-9.65);\draw[line width=.8pt] (6.3,-9.55) circle[radius=.19];\node[right,font=\footnotesize] at (6.499999999999999,-9.55) {smallest};
\end{tikzpicture}\\[1mm]{\footnotesize (a) How far each open statement was tested, and where a counterexample to each refuted one is known. Logarithmic scale.}\\[4mm]\begin{tabular}{@{}c@{\hspace{5mm}}c@{}}\begin{tikzpicture}[bf,bf small,x=0.47cm,y=0.47cm]
\path[use as bounding box] (-5.80,-5.80) rectangle (5.80,5.80);
\fill[bfblack,even odd rule] (0,0) circle[radius=5.55] (0,0) circle[radius=5.0];
\fill[bfblack] (-4.330,-2.500) -- (-2.133,-0.405) -- (-2.200,3.528) -- (-4.330,-2.500) -- cycle;
\fill[bfblack] (-4.330,-2.500) -- (-0.344,-4.394) -- (-1.614,-2.366) -- (-4.330,-2.500) -- cycle;
\fill[bfblack] (-2.133,-0.405) -- (0.764,-0.167) -- (-0.963,1.243) -- (-2.133,-0.405) -- cycle;
\fill[bfblack] (-2.133,-0.405) -- (-1.614,-2.366) -- (0.900,-2.302) -- (-2.133,-0.405) -- cycle;
\fill[bfblack] (-0.344,-4.394) -- (4.330,-2.500) -- (0.900,-2.302) -- (-0.344,-4.394) -- cycle;
\fill[bfblack] (4.330,-2.500) -- (2.799,2.992) -- (2.506,-0.651) -- (4.330,-2.500) -- cycle;
\fill[bfblack] (0.000,5.000) -- (-2.200,3.528) -- (2.799,2.992) -- (0.000,5.000) -- cycle;
\fill[bfblack] (-2.200,3.528) -- (-0.963,1.243) -- (0.610,2.065) -- (-2.200,3.528) -- cycle;
\fill[bfblack] (0.764,-0.167) -- (2.799,2.992) -- (0.610,2.065) -- (0.764,-0.167) -- cycle;
\fill[bfblack] (0.764,-0.167) -- (0.900,-2.302) -- (2.506,-0.651) -- (0.764,-0.167) -- cycle;
\draw[bf edge] (0,0) circle[radius=5.0];
\draw[bf edge] (-4.330,-2.500) -- (-2.133,-0.405);
\draw[bf edge] (-4.330,-2.500) -- (-1.614,-2.366);
\draw[bf edge] (-4.330,-2.500) -- (-0.344,-4.394);
\draw[bf edge] (-4.330,-2.500) -- (-2.200,3.528);
\draw[bf edge] (-2.133,-0.405) -- (-1.614,-2.366);
\draw[bf edge] (-2.133,-0.405) -- (-2.200,3.528);
\draw[bf edge] (-2.133,-0.405) -- (-0.963,1.243);
\draw[bf edge] (-2.133,-0.405) -- (0.764,-0.167);
\draw[bf edge] (-2.133,-0.405) -- (0.900,-2.302);
\draw[bf edge] (-1.614,-2.366) -- (-0.344,-4.394);
\draw[bf edge] (-1.614,-2.366) -- (0.900,-2.302);
\draw[bf edge] (-0.344,-4.394) -- (4.330,-2.500);
\draw[bf edge] (-0.344,-4.394) -- (0.900,-2.302);
\draw[bf edge] (4.330,-2.500) -- (0.900,-2.302);
\draw[bf edge] (4.330,-2.500) -- (2.506,-0.651);
\draw[bf edge] (4.330,-2.500) -- (2.799,2.992);
\draw[bf edge] (0.000,5.000) -- (-2.200,3.528);
\draw[bf edge] (0.000,5.000) -- (2.799,2.992);
\draw[bf edge] (-2.200,3.528) -- (-0.963,1.243);
\draw[bf edge] (-2.200,3.528) -- (2.799,2.992);
\draw[bf edge] (-2.200,3.528) -- (0.610,2.065);
\draw[bf edge] (-0.963,1.243) -- (0.764,-0.167);
\draw[bf edge] (-0.963,1.243) -- (0.610,2.065);
\draw[bf edge] (0.764,-0.167) -- (0.900,-2.302);
\draw[bf edge] (0.764,-0.167) -- (2.506,-0.651);
\draw[bf edge] (0.764,-0.167) -- (2.799,2.992);
\draw[bf edge] (0.764,-0.167) -- (0.610,2.065);
\draw[bf edge] (0.900,-2.302) -- (2.506,-0.651);
\draw[bf edge] (2.506,-0.651) -- (2.799,2.992);
\draw[bf edge] (2.799,2.992) -- (0.610,2.065);
\fill[af tint] (-2.133,-0.405) -- (0.900,-2.302) -- (0.764,-0.167) -- (-2.133,-0.405) -- cycle;
\begin{scope}\clip (-2.133,-0.405) -- (0.764,-0.167) -- (-0.963,1.243) -- (-2.133,-0.405) -- cycle;\path[bf s] (-0.963,1.243) -- ++(-125.37:0.709) arc[start angle=-125.37,end angle=-39.23,radius=0.709] -- cycle;\end{scope}
\begin{scope}\clip (-2.133,-0.405) -- (-1.614,-2.366) -- (0.900,-2.302) -- (-2.133,-0.405) -- cycle;\path[bf s] (-1.614,-2.366) -- ++(1.46:0.648) arc[start angle=1.46,end angle=104.82,radius=0.648] -- cycle;\end{scope}
\begin{scope}\clip (0.764,-0.167) -- (0.900,-2.302) -- (2.506,-0.651) -- (0.764,-0.167) -- cycle;\path[bf s] (2.506,-0.651) -- ++(164.47:0.723) arc[start angle=164.47,end angle=225.79,radius=0.723] -- cycle;\end{scope}
\begin{scope}\clip (-4.330,-2.500) -- (-2.133,-0.405) -- (-2.200,3.528) -- (-4.330,-2.500) -- cycle;\path[fill=black!45,draw=black!45] (-2.133,-0.405) -- ++(90.98:0.572) arc[start angle=90.98,end angle=223.64,radius=0.572] -- cycle;\end{scope}
\begin{scope}\clip (-4.330,-2.500) -- (-0.344,-4.394) -- (-1.614,-2.366) -- (-4.330,-2.500) -- cycle;\path[fill=black!45,draw=black!45] (-0.344,-4.394) -- ++(122.06:0.957) arc[start angle=122.06,end angle=154.58,radius=0.957] -- cycle;\end{scope}
\begin{scope}\clip (-0.344,-4.394) -- (4.330,-2.500) -- (0.900,-2.302) -- (-0.344,-4.394) -- cycle;\path[fill=black!45,draw=black!45] (0.900,-2.302) -- ++(-120.74:0.608) arc[start angle=-120.74,end angle=-3.30,radius=0.608] -- cycle;\end{scope}
\begin{scope}\clip (4.330,-2.500) -- (2.799,2.992) -- (2.506,-0.651) -- (4.330,-2.500) -- cycle;\path[fill=black!45,draw=black!45] (2.799,2.992) -- ++(-94.60:1.462) arc[start angle=-94.60,end angle=-74.42,radius=1.462] -- cycle;\end{scope}
\begin{scope}\clip (0.000,5.000) -- (-2.200,3.528) -- (2.799,2.992) -- (0.000,5.000) -- cycle;\path[fill=black!45,draw=black!45] (-2.200,3.528) -- ++(-6.12:1.042) arc[start angle=-6.12,end angle=33.79,radius=1.042] -- cycle;\end{scope}
\begin{scope}\clip (-2.200,3.528) -- (-0.963,1.243) -- (0.610,2.065) -- (-2.200,3.528) -- cycle;\path[fill=black!45,draw=black!45] (0.610,2.065) -- ++(152.50:0.710) arc[start angle=152.50,end angle=207.59,radius=0.710] -- cycle;\end{scope}
\begin{scope}\clip (0.764,-0.167) -- (2.799,2.992) -- (0.610,2.065) -- (0.764,-0.167) -- cycle;\path[fill=black!45,draw=black!45] (0.764,-0.167) -- ++(57.21:0.895) arc[start angle=57.21,end angle=93.95,radius=0.895] -- cycle;\end{scope}
\draw[bf Q,line width=2pt] (-2.133,-0.405) -- (0.764,-0.167);
\draw[bf Q,line width=2pt] (-2.133,-0.405) -- (0.900,-2.302);
\draw[bf Q,line width=2pt] (0.764,-0.167) -- (0.900,-2.302);
\draw[bf root] (-4.330,-2.500) circle[radius=1.4pt];
\fill[bf vertex] (-2.133,-0.405) circle[radius=.8pt];
\fill[bf vertex] (-1.614,-2.366) circle[radius=.8pt];
\fill[bf vertex] (-0.344,-4.394) circle[radius=.8pt];
\draw[bf root] (4.330,-2.500) circle[radius=1.4pt];
\draw[bf root] (0.000,5.000) circle[radius=1.4pt];
\fill[bf vertex] (-2.200,3.528) circle[radius=.8pt];
\fill[bf vertex] (-0.963,1.243) circle[radius=.8pt];
\fill[bf vertex] (0.764,-0.167) circle[radius=.8pt];
\fill[bf vertex] (0.900,-2.302) circle[radius=.8pt];
\fill[bf vertex] (2.506,-0.651) circle[radius=.8pt];
\fill[bf vertex] (2.799,2.992) circle[radius=.8pt];
\fill[bf vertex] (0.610,2.065) circle[radius=.8pt];
\end{tikzpicture} & \begin{minipage}[b][5.4cm][c]{7.6cm}\footnotesize (b) The mechanism behind the failures at the extremes, on the state $s_{9}$ of the running example. The three black faces round the hatched white face all put their wedge (blue) at the corner away from it, so their opposite edges close a triangle (heavy dashes). The 12-vertex counterexample on page 16 fails through a face of this kind.\end{minipage}\end{tabular}\end{tabular}%
\end{document}
